\subsection{Case Study: A Published SMT Assembly Line}
\label{subsec:smt_case_study}

The preceding subsections evaluated FactoryFlow against reference systems constructed for
this work. This case study evaluates it against a published system, so that the structure,
the parameters and the reported KPIs all originate outside this paper.

\noindent\textbf{The system.} The case system is an automated Surface Mount Technology
(SMT) printed circuit board assembly line operating at a manufacturing facility in Mysuru,
India, and composed of legacy machines \cite{KaranjkarSMT2018}. It was instrumented with
energy meters, vibration sensors and infrared proximity sensors, and machine states and
state-wise power draw were inferred from three months of recorded data. Boards arrive in
stacks of sixteen and pass through five machines in sequence: a line loader, a screen
printer, two pick-and-place machines and a reflow oven, linked by single-board transfer
positions and two non-accumulating conveyor belts. One human operator serves the whole
line, refilling the printer's consumables and replacing component reels on the placement
machines. The published study provides both the model and its result files, so the
parameters are exact rather than transcribed from a summary. Its reported values are
outputs of that model: the study states agreement with the real line in prose but does not
publish the measured values, so the comparisons below are against a published reference
model of a real line rather than against measured factory throughput.

\noindent\textbf{Experimental methodology.} A description of the line was supplied to the
FactoryFlow model generator, which produced an executable FactorySimPy model. The
generated model was executed unmodified to the same terminating condition as the
reference, namely until 1024 boards had been completed, and its KPIs were compared with
the published values. The system description is an account of the line of the kind a
production engineer would give, naming the components and stating their timings and power
draw. Because model generation is not deterministic, ten candidates were generated from
the same system description with Gemini~2.5~Pro.

\noindent\textbf{Model size and parameters.} The line comprises eight flow nodes, one
operator and seven edges: a source, five machines, a pass-through module and a sink,
connected by five single-board transfer positions and two conveyor belts. Its behaviour is
governed by five processing delays, one cleaning delay and the interval at which cleaning
occurs, three service intervals and the operator task duration, two belt parameters, three
oven parameters and the stack size. Five of these are processing delays obtainable from
telemetry; the remainder are structural or configuration attributes taken from the
description.

Two components exercise capabilities beyond a machine with a processing delay. The refills
and reel changes are performed by a single shared operator, so a machine whose service
falls due stalls until the operator is free; the operator is a resource that machines
contend for rather than an additional delay folded into processing time. The reflow oven
holds ten boards simultaneously and must reach temperature before it accepts any board,
and is expressed through component capacity and setup time rather than as a chain of
stages.

\noindent\textbf{Structural agreement.} Every generated candidate was compared with the
reference graph by component identifier, component type and the endpoints of each edge.
Six of the ten candidates executed, and all six reproduced the reference graph exactly:
the same nine nodes, the same seven edges and the same endpoints, including the
recognition that boards arriving in stacks require a batch-releasing source rather than a
per-item one. Structure was therefore recovered reliably from the system description,
without component or connection errors.

\noindent\textbf{KPI agreement.} All six executing candidates matched the published line's
reported throughput, maximum cycle time, energy per board and reflow-oven warm-up fraction
to within 0.01\%. Table~\ref{tab:smt_baseline} gives the values for one candidate; the
other five agree with the published values to the same precision.

\begin{table}[!t]
\centering
\caption{SMT line: generated digital twin against the published reference model.}
\label{tab:smt_baseline}
\footnotesize
\setlength{\tabcolsep}{4pt}
\begin{tabular}{@{}lrrr@{}}
\toprule
KPI & Published & DT & Diff. \\
\midrule
Throughput (boards/h)      & 40.3549 & 40.3547 & $0.00\%$ \\
Maximum cycle time (h)     & 1.6193  & 1.6194  & $+0.01\%$ \\
Energy per board (kJ)      & 2427.29 & 2427.65 & $+0.01\%$ \\
Oven warm-up fraction (\%) & 3.9409  & 3.9409  & $0.00\%$ \\
\bottomrule
\end{tabular}
\end{table}

Energy follows from the per-state power ratings the
generated model declares on its components together with the time each component spends in
each state, so it is an output of the generated model rather than a separate calculation.
The operator was occupied for 3.94\% of the run and no service request queued, so operator
contention does not affect this configuration; this is measured rather than assumed, and
would not have been visible had the service time been folded into the processing delays.

Mean cycle time is excluded because it is not comparable. The reference creates a stack of
sixteen boards and timestamps all sixteen at the instant the stack is created, so a board
accumulates waiting time before it enters the line. Instrumenting per-board residence
times shows that changing how many stacks wait at the head of the line changes only this
pre-line interval; the time a board spends inside the components is unchanged. The maximum
cycle time is set by the initial oven warm-up and is therefore unaffected by the
convention.

\noindent\textbf{Assumption review.} The system description leaves one behavioural
question open: it states that the oven must reach temperature before it accepts boards,
and gives a power rating for the oven when switched off, but does not say whether the oven
ever switches off in normal operation. The answer determines whether the oven warms up
once at startup or before every batch, and so governs throughput, cycle time and energy.
The generator records the assumptions it made alongside the model it produced, and every
candidate that executed stated this one explicitly: the oven warms up once and then
remains at temperature, as in the published model.

The choice was therefore visible before any model was run. Had a candidate resolved it the
other way, an expert familiar with the line would correct the stated assumption directly,
and the revised description or assumption set would yield a model with the intended
behaviour. This is the acceptance and revision mechanism of Stage~C: ambiguity in the
description is surfaced as a stated assumption for confirmation instead of being silently
resolved. The case is consistent with the division of labour the framework assumes:
structural inference from a description, parameter values from telemetry where it is
available (here they were stated in the description), and expert confirmation of the
assumptions that remain.

\noindent\textbf{Scope.} The published study evaluates the line in three
configurations: the line as installed, and two variants in which a buffering
module accumulates boards so that the reflow oven can be switched off between
batches. This case study considers only the first. Its purpose is not to repeat
the energy study but to show that FactoryFlow, and the FactorySimPy library it
composes, extend beyond constructed reference systems to a line in industrial
operation, with a shared human operator, a batch-fed source, and a multi-board
oven that must reach temperature before it accepts any board.

\noindent\textbf{Limitations.} The comparison targets are the published model's reported
values rather than measured factory KPIs, for the reason given above, and mean cycle time
is not comparable because of the timestamping convention.

Candidate generation is not deterministic. Four of the ten candidates failed to execute
because they used a library interface that does not exist: two imported an absent module,
and two declared power ratings for a machine state that FactorySimPy does not define.
These are not resolvable by revising the description: they are caught by the execution
check of Stage~C and handled by regeneration, and none produced a structurally incorrect
model.
